% =====================================================================================
% Manuscript changes for the red items in the reply letter (round 2), marked in ORANGE.
% Paste each block into manuscript.tex at the place indicated. \textcolor{orange}{...}
% needs xcolor, which the manuscript already loads (it uses blue for the earlier changes).
%
% Macros used: \vek, \grad, \p as in the original preamble, and \bmag for the rescaled
% magnetic field b. If your preamble uses a different name for b, replace \bmag below.
% Labels used: fig:Gaussian_field (the Gaussian field, now Fig. 4), eq:efflag (Eq. 1).
% =====================================================================================


% -------------------------------------------------------------------------------------
% (A) Reply point 3 (red: "dark refers to small positive and light to large positive?")
%     Caption of Fig. 5 (sharp domain wall, chiral limit, b_0 = 5, 10 x 10 domain).
%     Insert right after the sentence
%       "The baryon number density vanishes along the domain wall localized at $x=0$
%        where the field changes sign."
% -------------------------------------------------------------------------------------

\textcolor{orange}{Unlike in Fig.~\ref{fig:Gaussian_field}, the baryon number density is nonnegative everywhere; it reaches its largest value, $\bmag_0^2$, on the boundaries $z=\pm L_z/2$. The grayscale therefore runs from dark where the baryon number density vanishes, along the domain wall, to light where it is largest, near the top and bottom of the domain.}


% -------------------------------------------------------------------------------------
% (B) Reply point 4 (red: "What is the actual scale of variation of magnetic fields for
%     the magneto-rotational instabilities?")
%     Introduction, scales paragraph. Insert right after the sentence ending
%       "... would likely destroy any CSL-type order."
%     New references: Siegel2013MRI, Guilet2017MRI (entries in new_references_round2.bib).
% -------------------------------------------------------------------------------------

\textcolor{orange}{Here ``small-scale'' means small compared to the size of the star, not compared to the scale of QCD. For instance, the wavelength of the fastest-growing mode of the magnetorotational instability, $\lambda_\text{MRI}\simeq2\pi v_\text{A}/\Omega$, where $v_\text{A}=B/\sqrt{4\pi\rho}$ is the Alfv\'en velocity (in Gaussian units), $\rho$ the mass density and $\Omega$ the angular velocity, ranges from about a meter for $B\sim10^{13}~\mathrm{G}$ to kilometers for $B\gtrsim10^{15}~\mathrm{G}$, assuming $\rho\sim10^{13}$--$10^{14}~\mathrm{g}/\mathrm{cm}^3$ and $\Omega\sim10^3$--$10^4~\mathrm{s}^{-1}$, as appropriate for merger remnants~\cite{Siegel2013MRI,Guilet2017MRI}. This exceeds the scale $1/m_\pi$ by at least fifteen orders of magnitude.}


% -------------------------------------------------------------------------------------
% (C) Reply point 5 (red: "To be implemented.")
%     Sec. 5 (Summary and outlook). Insert as a new paragraph right BEFORE the paragraph
%       "There are at least two natural avenues for follow-up research in this direction."
%     Uses only references already in the bibliography (cited in the introduction for the
%     excitation spectrum): Yamamoto:2015maz = arXiv:1512.05668, Brauner:2017mui =
%     Brauner & Kadam, JHEP 03 (2017) 015. Please check the keys point to these papers.
% -------------------------------------------------------------------------------------

\textcolor{orange}{Throughout this paper, we have assumed, in line with most of the literature on the CSL phase of quark matter, that the magnetic field is fixed and does not experience any feedback from the neutral pion condensate. The sources of the magnetic field are not included in our Lagrangian~\eqref{eq:efflag}, and may consist of other types of matter than pions. One therefore cannot expect the magnetic field to satisfy the Maxwell equations that would follow by adding the kinetic term for the electromagnetic field to the Lagrangian~\eqref{eq:efflag}. It is possible to treat the electromagnetic field as a dynamical degree of freedom coupled to the pions, which leads to interesting and highly nontrivial effects well beyond the scope of the present paper. Namely, in analogy with axion electrodynamics, the anomalous term in the Lagrangian then mixes neutral pions and photons in the presence of a background magnetic field, which has dramatic consequences for the photon spectrum as well as for the low-temperature thermodynamics of the system~\cite{Yamamoto:2015maz,Brauner:2017mui}.}


% =====================================================================================
% OPTIONAL (not marked red, but still in the current text): language fixes in the blue
% scales paragraph of the introduction. Apply if you agree; they stay blue.
%   - "$\frac{1}{m_{\pi}}$" (twice)        ->  "$1/m_\pi$"
%   - "a few $fm$"                         ->  "a few~fm"
%   - "at RHIC energies"                   ->  "at the top RHIC energy"
%   - "Another parameter, the baryon chemical potential is small ... and the temperature"
%       -> "Moreover, the baryon chemical potential is small ..., and the temperature"
%   - last sentence has no main verb:
%       "The inhomogeneities considered in this paper, which vary on a scale comparable to
%        $1/m_\pi$ and are therefore far from ..."
%       -> "The inhomogeneities considered in this paper, which vary on a scale comparable
%           to $1/m_\pi$, are therefore far from ..."
%   - Fig. 4 caption and Eq. (30) still use $\psi_0$, $\psi$ for the Gaussian stream
%     function, which clashes with the scalar potential $\psi[\vek\bmag]$ of Eq. (17).
% =====================================================================================
